% ======================================================================
% SECTION 7: HR 9507 - Cancer Patient Health Data Self-Custody and
% Trial-Selection Act of 2026
% Revision of HHS HIPAA Right-to-Access 2025 and ONC Cures Act Final
% Rule.
%
% This file consolidates instructions previously housed in the prior
% template's Patient Priority Framework dimension 2.4 (data and
% biospecimen governance), Implementation/Metrics (Domain 1 access
% subset on real-time data release), Discussion (11.1.a Cures Act
% information-blocking extension and 11.2.c Virchow pathology
% evidence and 11.3.d HHS HIPAA opening), Limitations (12.1.g HR
% 4507 limitation), and Conclusions (13.1 seventh list item) into
% the body of HR 9507.
% ======================================================================

[CONSOLIDATED PATIENT PRIORITY FRAMEWORK DIMENSION 2.4 (formerly
Patient Priority Framework subsection 2.4). Data and biospecimen
governance: HIPAA, broad consent, information-blocking penalties.
Cite \cite{hhs_right_to_access_2025},
\cite{onc_cures_act_final_rule_page}, and
\cite{federal_register_info_blocking_disincentives_2024}.]

[CONSOLIDATED INTRODUCTION TIE-IN (formerly Introduction block 2
items 2.a and 2.h). 21st Century Cures Act, PL 114-255 (2016):
interoperability, information-blocking, real-world evidence, and
the non-device clinical decision support software carve-out. HHS
HIPAA Right-to-Access (2025), ASTP/ONC Patient Portals Brief 2024,
and ASTP/ONC Hospital APIs Brief 2024. Information-blocking
penalties exist (\cite{federal_register_info_blocking_disincentives_2024})
but patient-initiated self-custody of trial-eligibility data is not
yet guaranteed. HR 9507 closes that gap. Cite
\cite{rank1_cures_act_2016},
\cite{onc_cures_act_final_rule_page},
\cite{hhs_right_to_access_2025},
\cite{astp_onc_patient_portals_apps_2024}, and
\cite{astp_onc_hospital_apis_2024}.]

[PROCESSING INSTRUCTION (HR 9507 - Cancer Patient Health Data Self-
Custody and Trial-Selection Act of 2026, populate to 9 to 12
paragraphs structured into the five legislative-act subsections):

The bill must explicitly state in its first paragraph that it is a
"Revision of HHS HIPAA Right-to-Access (2025) and the ONC Cures Act
Final Rule" so that readers can trace the chain of prior authority.
Cite \cite{hhs_right_to_access_2025},
\cite{onc_cures_act_final_rule_page}, and
\cite{federal_register_info_blocking_disincentives_2024}.

Subsection 7.1 - Findings and Declarations: Read in full
- patients/paper/Deep-Research-3/part1_legal_baseline.md (the
  Cures Act / HIPAA data-access rights baseline)
- patients/paper/Deep-Research-3/part3_metrics_guardrails.md (the
  four-domain dashboard, in particular Domain 1 access and Domain
  4 fairness and technical quality)
- new-trial/site/03-legislation-data-transparency/physical_ai_data_transparency.tex
  and national-platform/new_trial_psl/03_sb892_data_protection.tex
  (the prior CA SB 892 Clinical Data Protection and Transparency
  Act baseline that this national bill revises into federal
  coverage)

Compose 5 to 7 enumerated findings citing
\cite{astp_onc_patient_portals_apps_2024},
\cite{astp_onc_hospital_apis_2024},
\cite{vorontsov2024virchow} (for the pathology data that needs
same-day release for trial-eligibility gating), and the prior CA
SB 892 framework being revised into federal coverage.

Subsection 7.2 - Definitions:
- "Patient health data self-custody" means the patient's right to
  hold their own complete oncology medical record, in machine-
  readable FHIR format, on a patient-controlled device or cloud
  endpoint, with full HIPAA portability across providers, payers,
  and trial sponsors.
- "Trial-eligibility data" means the subset of the patient's
  medical record that determines eligibility for currently
  recruiting oncology trials, including pathology slides and
  reports, genomic sequencing reports, prior treatment history,
  ECOG/PD-L1/TMB scores, and current medications. Reference
  Virchow pathology benchmark data at \cite{vorontsov2024virchow}
  and the patient-journey stage_03_digital_twin.py module for the
  technical baseline.
- "Same-day release" means pathology and genomic results released
  to the patient's self-custody endpoint within the same calendar
  day the result is finalized, drawing on Domain 1 of the four-
  domain dashboard at
  patients/paper/Deep-Research-3/part3_metrics_guardrails.md.

Subsection 7.3 - Patient Health Data Self-Custody Rights (operative
core):
- Right to receive the complete medical record in machine-readable
  FHIR format on demand, at zero cost to the patient.
- Right to require same-day release of pathology and genomic
  results to the patient's self-custody endpoint, drawing on the
  ASTP/ONC hospital APIs framework at
  \cite{astp_onc_hospital_apis_2024}.
- Right to grant time-limited, scope-limited access to specific
  trial sponsors or AI matching systems (TrialGPT, PRISM,
  TrialMatchAI), with revocation taking effect within 1 hour. Cite
  \cite{jin2024trialgpt}, \cite{gupta2024prism}, and
  \cite{abdallah2026trialmatchai}.
- Right to a portable trial-history record that follows the
  patient across providers, payers, and states; reference the ICH
  E6(R3) adaption at
  national-platform/ich_e6r3_adapt/03_data_governance.tex for the
  data-governance baseline. Cite \cite{kawchak_2026_18973368}.
- Right to information-blocking penalties under
  \cite{federal_register_info_blocking_disincentives_2024} where
  any provider or payer denies the patient's self-custody
  request.
- Right to a permanent, exportable copy of the complete oncology
  medical record at end of trial enrollment, with no provider
  retention requirement that overrides the patient's self-custody
  rights.

Subsection 7.4 - Implementation: 12-month deadline for ALL CMS-
billable EHR vendors to publish a self-custody export endpoint;
24-month deadline for state-level harmonization; FDA, ONC, OCR,
and CMS roles; preemption rule preserves stronger state
protections (e.g., California's prior data-protection statute
baseline).

Subsection 7.5 - Reporting and Enforcement: Annual federal report
on per-state, per-payer self-custody uptake; civil penalties under
HIPAA enforcement and information-blocking disincentives; CMS
reimbursement penalty for sites that fail to expose the self-
custody endpoint.]

[CONSOLIDATED IMPLEMENTATION-METRICS DOMAIN 1 ACCESS SUBSET
(formerly Implementation Metrics subsection 10.2 Domain 1 subset on
real-time data release). The federal Patient Control Dashboard
must publish, for HR 9507, the following access metrics
specifically tied to self-custody: same-day pathology release rate,
same-day genomic release rate, end-to-end FHIR export latency,
patient-initiated trial sponsor access grants per quarter, and
patient-initiated revocations within 1 hour. Cite
\cite{hhs_right_to_access_2025} and \cite{astp_onc_hospital_apis_2024}.]

[CONSOLIDATED DISCUSSION COMPARISON (formerly Discussion items
11.1.a Cures Act information-blocking extension, 11.2.c Virchow
pathology evidence, and 11.3.d HHS HIPAA opening). HR 9507 extends
the prior 21st Century Cures Act information-blocking framework
into a positive self-custody right. The Virchow foundation-model
performance across pan-cancer and rare-cancer detection supports
HR 9507's same-day pathology release requirement. The HHS HIPAA
Right-to-Access 2025 framework is the regulatory opening that HR
9507 builds on. Cite \cite{rank1_cures_act_2016},
\cite{onc_cures_act_final_rule_page},
\cite{vorontsov2024virchow}, and
\cite{hhs_right_to_access_2025}.]

[CONSOLIDATED LIMITATIONS PARAGRAPH (formerly Limitations item
12.1.g). HR 9507 assumes EHR vendors will support a self-custody
export endpoint at zero cost to the patient within 12 months. The
Cures Act information-blocking enforcement record is mixed. Cite
\cite{federal_register_info_blocking_disincentives_2024}.]

[CONSOLIDATED CONCLUSIONS RESTATEMENT (formerly Conclusions item
13.1 seventh list item). HR 9507 replaces the prior access-only
HIPAA framework with a positive self-custody right and a
trial-selection authorization mechanism. Cite
\cite{kawchak_2026_19119939} and \cite{kawchak_2026_19994945}.]

Close the bill (1 paragraph after Subsection 7.5) with two
patient-control framings:

(i) FOR INDIVIDUAL PATIENTS: A 58-year-old female NSCLC patient
    matching PAT-2026-0042 receives her PD-L1 65 percent and TMB
    14 mut/Mb scores released to her self-custody endpoint on the
    same calendar day, grants 30-day time-limited access to
    TrialGPT and TrialMatchAI for trial discovery (HR 9501), and
    revokes that access within 1 hour after self-booking her
    preferred trial.

(ii) FOR BROAD ADOPTION: All CMS-billable EHRs must expose a
     self-custody export endpoint within 12 months; all U.S.
     cancer patients gain portable trial-eligibility data within
     24 months. The bill positions the United States as Number 1
     in the world for patient health data self-custody, replacing
     the prior access-only HIPAA framework with a positive
     self-custody right. Cite \cite{kawchak_2026_19119939} and
     \cite{kawchak_2026_19994945}.
